\notechapter{N-20220911}{Analytic Geometry through Normals, Line Pencils, and Quadratic Forms}

\section{A line is an affine level set, not a collection of unrelated formulas}

A line in the plane can be written
\[
  H=\{x\in\mathbb R^2:n\cdot x=c\},
  \qquad n\ne0.
\]
The vector \(n=(a,b)\) is normal to the line, and a direction vector is
\[
  d=(-b,a).
\]
If \(p\in H\), every point of the line has the parametric form
\[
  x=p+td,
  \qquad t\in\mathbb R.
\]
Conversely, eliminating \(t\) gives
\[
  a(x_1-p_1)+b(x_2-p_2)=0.
\]
This single description includes vertical lines, where slope notation fails.

The familiar equations are coordinate choices for the same object.  If \(b\ne0\),
\[
  ax+by=c
  \quad\Longleftrightarrow\quad
  y=-\frac abx+\frac cb.
\]
If a line passes through \(p_1,p_2\), its direction is \(p_2-p_1\).  If it meets the coordinate axes away from the origin, dividing by its intercepts produces
\[
  \frac xa+\frac yb=1.
\]
None of these forms is universally best; the normal form survives every orientation, while slope form is efficient only in the nonvertical chart.

In homogeneous coordinates, the line is represented by the coefficient triple
\[
  \ell=[a:b:-c],
\]
and incidence with a projective point \(X=[x:y:z]\) is the bilinear equation
\[
  ax+by-cz=0.
\]
The homogeneous form will later make points at infinity and polar lines ordinary parts of the calculation.

\section{Angles and Euclidean distance come from the normal vector}

Two lines with normals \(n_1,n_2\) are parallel exactly when the normals are proportional.  They are perpendicular exactly when
\[
  n_1\cdot n_2=0,
\]
because the direction of each line is a quarter-turn of its normal.

For
\[
  H=\{x:n\cdot x=c\}
\]
and a point \(x_0\), the signed displacement in the normal direction is
\[
  \frac{n\cdot x_0-c}{\|n\|_2}.
\]
Hence
\[
  \boxed{
  \operatorname{dist}_2(x_0,H)
  =\frac{|n\cdot x_0-c|}{\|n\|_2}.}
\]
To derive it, let
\[
  p=x_0-\frac{n\cdot x_0-c}{\|n\|_2^2}n.
\]
Then \(n\cdot p=c\), so \(p\in H\), and \(x_0-p\) is normal to the line.  Any other point of \(H\) differs from \(p\) by a tangent vector, so the Pythagorean theorem makes \(p\) the closest point.

For parallel lines
\[
  n\cdot x=c_1,
  \qquad
  n\cdot x=c_2,
\]
the same formula gives
\[
  \operatorname{dist}_2(H_1,H_2)
  =\frac{|c_1-c_2|}{\|n\|_2}.
\]
The denominator is not a memorized normalization constant; it compensates for rescaling the defining equation.

\section{Distance to a hyperplane is governed by the dual norm}

The Euclidean formula has a normed-space version.  Let \(\|\cdot\|\) be any norm on \(\mathbb R^n\), and let \(\varphi\) be a nonzero linear functional.  Its dual norm is
\[
  \|\varphi\|_*
  =\sup_{\|u\|\le1}|\varphi(u)|.
\]
For the hyperplane
\[
  H=\{x:\varphi(x)=c\},
\]
one has
\[
  \boxed{
  \operatorname{dist}(x_0,H)
  =\frac{|\varphi(x_0)-c|}{\|\varphi\|_*}.}
\]
Indeed, for every \(x\in H\),
\[
  |\varphi(x_0)-c|
  =|\varphi(x_0-x)|
  \le\|\varphi\|_*\|x_0-x\|,
\]
which gives the lower bound.  In finite dimensions, the supremum defining the dual norm is attained by a unit vector \(u\); moving from \(x_0\) in the signed \(u\)-direction reaches equality.

For \(\varphi(x,y)=ax+by\), the dual of \(\ell^1\) is \(\ell^\infty\), so
\[
  \operatorname{dist}_1((x_0,y_0),H)
  =\frac{|ax_0+by_0-c|}{\max(|a|,|b|)}.
\]
The dual of \(\ell^\infty\) is \(\ell^1\), giving
\[
  \operatorname{dist}_\infty((x_0,y_0),H)
  =\frac{|ax_0+by_0-c|}{|a|+|b|}.
\]
The familiar square root in the Euclidean formula is therefore a special case of dual-norm geometry.

\section{Circles are quadratic level sets, and subtraction reveals their common geometry}

A circle with center \(m=(a,b)\) and radius \(r\) is
\[
  \|x-m\|_2^2=r^2.
\]
Expanding gives
\[
  x^2+y^2+Dx+Ey+F=0,
\]
while completing squares recovers
\[
  m=\left(-\frac D2,-\frac E2\right),
\]
\[
  r^2=\frac{D^2+E^2-4F}{4}.
\]
The sign of this final quantity distinguishes a real circle, a point circle, and no real locus.

For a point \(p\), define its power with respect to the circle by
\[
  \operatorname{Pow}_C(p)=\|p-m\|^2-r^2.
\]
If a line through \(p\) meets the circle at \(a,b\), then
\[
  \operatorname{Pow}_C(p)=\overline{pa}\,\overline{pb}
\]
using directed lengths.  This follows by substituting the line parametrization \(p+td\) into the circle equation: the product of the two parameter roots is the constant term divided by the leading coefficient.

For two circles, the locus of equal powers is their radical axis.  Subtracting
\[
  C_1:x^2+y^2-4x=0
\]
and
\[
  C_2:x^2+y^2-2y-3=0
\]
cancels the quadratic terms and gives the line
\[
  -4x+2y+3=0.
\]
The radical axis is linear because it compares two quadratic functions with the same leading form.

\section{A fractional range problem is a pencil of lines through one point}

Consider
\[
  \mu=\frac{x}{y-3}
\]
under the constraint
\[
  |x|+|y|\le1.
\]
For fixed \(\mu\), rewrite the equation as
\[
  x=\mu(y-3).
\]
This is a line through
\[
  P=(0,3).
\]
The range problem asks which lines in the pencil through \(P\) meet the diamond
\[
  K=\{(x,y):|x|+|y|\le1\}.
\]
Because \(K\) is convex and \(P\notin K\), the extreme parameters occur at supporting lines from \(P\) to \(K\).

The support function of the \(\ell^1\)-unit ball is
\[
  h_K(n)=\max_{x\in K}n\cdot x=\|n\|_\infty.
\]
The line can be written
\[
  (1,-\mu)\cdot(x,y)=-3\mu.
\]
It meets \(K\) precisely when
\[
  |3\mu|\le h_K(1,-\mu)=\max(1,|\mu|).
\]
If \(|\mu|>1\), this is impossible.  If \(|\mu|\le1\), it becomes
\[
  3|\mu|\le1.
\]
Therefore
\[
  \boxed{
  \mu\in\left[-\frac13,\frac13\right].}
\]
The endpoints correspond to the supporting lines through the vertices \((1,0)\) and \((-1,0)\).  Convex duality has not replaced the elementary picture; it has explained why the endpoints are support data.

\section{Ratios on a graph are radial coordinates}

For a point \((x,f(x))\) with \(x\ne0\), the ratio
\[
  \frac{f(x)}x
\]
is the slope of the ray from the origin to the graph point.  Counting solutions of
\[
  \frac{f(x)}x=m
\]
is therefore the same as counting intersections between the graph and the line
\[
  y=mx.
\]
A rational expression has become a radial-projection problem.

This viewpoint also identifies tangency thresholds.  As \(m\) varies, the number of intersections can change only when the line becomes tangent to the graph or passes through a singular point.  Algebraically, solving
\[
  f(x)=mx
\]
together with
\[
  f'(x)=m
\]
detects those critical slopes.  The graphical and derivative descriptions are the same incidence event viewed globally and locally.

\section{A perpendicular moving-line problem is a circle problem}

Let fixed points be
\[
  A=(-3,0),
  \qquad
  B=(1,6).
\]
Suppose lines through \(A\) and \(B\) meet at \(P\) and are perpendicular.  Then
\[
  \angle APB=90^\circ,
\]
so Thales' theorem places \(P\) on the circle with diameter \(AB\).

The right triangle gives
\[
  PA^2+PB^2=AB^2.
\]
Hence
\[
  PA\,PB
  \le\frac{PA^2+PB^2}{2}
  =\frac{AB^2}{2}.
\]
Since
\[
  AB^2=(1+3)^2+6^2=52,
\]
one obtains
\[
  \boxed{
  PA\,PB\le26.}
\]
Equality occurs when \(PA=PB\), namely at the intersections of the diameter circle with the perpendicular bisector of \(AB\).

The coordinate version would introduce two slopes and a quadratic constraint.  The circle formulation identifies the locus first, after which the optimization is a one-line inequality.

\section{A general conic is an affine quadratic form}

Write an affine conic as
\[
  F(x)=x^TQx+2\ell^Tx+f=0,
  \qquad Q=Q^T.
\]
An orthogonal change of coordinates diagonalizes \(Q\), removing the cross term without changing Euclidean angles.  If \(Q\) is invertible, translating by
\[
  x=u-Q^{-1}\ell
\]
removes the linear term.  Rotation handles directions; translation handles the center.

Consider
\[
  xy-x-y=0.
\]
Set
\[
  u=\frac{x+y}{\sqrt2},
  \qquad
  v=\frac{x-y}{\sqrt2}.
\]
Then
\[
  xy=\frac{u^2-v^2}{2},
  \qquad
  x+y=\sqrt2\,u.
\]
The equation becomes
\[
  u^2-v^2-2\sqrt2\,u=0,
\]
or
\[
  \boxed{
  (u-\sqrt2)^2-v^2=2.}
\]
The original rectangular-looking equation is a hyperbola whose principal axes are the eigenvectors of the quadratic matrix
\[
  Q=
  \begin{pmatrix}
    0&1/2\\
    1/2&0
  \end{pmatrix}.
\]

\section{Rank and the line at infinity distinguish conic types}

Homogenizing a conic gives
\[
  X^TAX=0,
  \qquad X=(x,y,z)^T,
\]
for a symmetric \(3\times3\) matrix \(A\).  The conic is degenerate when \(A\) loses rank.  For example,
\[
  xy=0
\]
is the union of two lines and has a rank-deficient homogeneous matrix.

The intersection with the line at infinity \(z=0\) records affine directions.

\begin{itemize}
  \item An ellipse
  \[
    \frac{x^2}{a^2}+\frac{y^2}{b^2}=z^2
  \]
  has no real point at infinity.
  \item A parabola
  \[
    y^2=2pxz
  \]
  meets the line at infinity at one double point \([1:0:0]\).
  \item A hyperbola
  \[
    \frac{x^2}{a^2}-\frac{y^2}{b^2}=z^2
  \]
  has two real points at infinity, corresponding to its asymptotic directions.
\end{itemize}

For
\[
  xy-x-y=0,
\]
the projective closure is
\[
  XY-XZ-YZ=0.
\]
At infinity,
\[
  Z=0
  \quad\Longrightarrow\quad
  XY=0,
\]
so the two points are
\[
  [1:0:0],
  \qquad
  [0:1:0].
\]
They are the projective endpoints of the two asymptote directions revealed by the standard hyperbola form.

\section{Polarity turns the conic matrix into its tangent geometry}

Let a nondegenerate projective conic be
\[
  F(X)=X^TAX=0,
  \qquad A=A^T.
\]
For a point \(p\), its polar line is
\[
  p^TAX=0.
\]
If \(p\) lies on the conic, this polar is the tangent line.  Indeed, the differential is
\[
  dF_p(H)=2p^TAH.
\]
The tangent condition at \(p\) is
\[
  p^TA(X-p)=0.
\]
Since \(p^TAp=0\), this reduces to
\[
  p^TAX=0.
\]

For the conic
\[
  XY-XZ-YZ=0,
\]
use the symmetric matrix
\[
  A=
  \begin{pmatrix}
    0&1&-1\\
    1&0&-1\\
    -1&-1&0
  \end{pmatrix},
\]
which represents twice the equation.  The affine point \((2,2)\) corresponds to
\[
  p=[2:2:1]
\]
and lies on the conic.  Since
\[
  Ap=
  \begin{pmatrix}1\\1\\-4\end{pmatrix},
\]
its polar, hence its tangent, is
\[
  X+Y-4Z=0.
\]
In the affine chart \(Z=1\),
\[
  \boxed{x+y=4.}
\]
The same answer follows from the gradient of \(xy-x-y\) at \((2,2)\), but the polar calculation remains valid in homogeneous coordinates and transforms naturally with the conic under projective changes of coordinates.